\documentclass[a4paper,12pt]{ctexart}
\usepackage[margin=2.54cm]{geometry}
\usepackage{amsmath,newtxmath,newtxtext,bm}
\usepackage{minted}
\usepackage[svgnames]{xcolor}
\setminted{
  breaklines,
  breakanywhere,
  autogobble,
  fontsize=\small
}
% \usepackage{indentfirst}
\setlength{\parindent}{2em}
\setlength{\parskip}{0pt}
\linespread{1.25}
\numberwithin{equation}{section}
\pagestyle{plain}

\usepackage[
  colorlinks=true,
  linkcolor=blue,
  citecolor=blue,
  urlcolor=blue]{hyperref}

\usepackage[
  backend=biber,
  style=numeric,
  sorting=none]{biblatex}
\addbibresource{cute_algebra.bib}

\begin{document}
\section{CUTE layout代数}
\subsection{一些基本的layout操作}
A和B为CUTE layout，那么A与B的logical division为：$A \oslash B = A \circ (B,B^c)$，其中$B^c=\text{comp}(B,\text{size}(A))$。
logical product定义为：$A \otimes B = (A,A^c \circ B)$，其中$A^c=\text{comp}(A,\text{size}(A)*\text{cosize}(B))$。
更多的关于layout的基本操作可以查询参考文献\cite{carlisle2026categoricalfoundationscutelayouts}。
\subsection{sgemm\_1.cu}
在cute的第一个例子sgemm\_1.cu中，
\begin{minted}[bgcolor=Beige,linenos]{cpp}
  // Represent the full tensors
  // (M,K)
  Tensor mA = make_tensor(make_gmem_ptr(A), select<0, 2>(shape_MNK), dA);
  // (N,K)
  Tensor mB = make_tensor(make_gmem_ptr(B), select<1, 2>(shape_MNK), dB);
  // (M,N)
  Tensor mC = make_tensor(make_gmem_ptr(C), select<0, 1>(shape_MNK), dC);
\end{minted}
上述代码片段变量的layout分别为：
\begin{equation}
  \begin{aligned}
    \text{mA} = (5120,4096):(1,5120) \\
    \text{mB} = (5120,4096):(1,5120) \\
    \text{mC} = (5120,5120):(1,5120) \\
  \end{aligned}
\end{equation}

\printbibliography[title={参考文献}]
\end{document}
